\par\Needspace{6\baselineskip}\medskip
\SourceLeaf{158}\pdfbookmark[2]{113.}{article-113}
\begin{center}113.\end{center}
Ne formularum prolixitate nimis obruamur, sequentibus abbreviationibus uti placet. Primo designamus quantitatem
\[
\operatorname{tang}\beta\sin(\alpha''-\alpha')+\operatorname{tang}\beta'\sin(\alpha-\alpha'')+\operatorname{tang}\beta''\sin(\alpha'-\alpha)
\]
per $(0.1.2)$: si in expressione illa pro longitudine et latitudine loco cuivis geocentrico respondentibus substituuntur longitudo et latitudo cuilibet trium locorum heliocentricorum terrae respondentes, in signo $(0.1.2)$ numerum illi respondentem cum numero romano eo commutamus, qui posteriori respondet. Ita e. g. character $(0.1.\mathrm{I})$ exprimet quantitatem
\[
\operatorname{tang}\beta\sin(L'-\alpha')+\operatorname{tang}\beta'\sin(\alpha-L')+\operatorname{tang}B'\sin(\alpha'-\alpha)
\]
nec non character $(0.\mathrm{O}.2)$ hanc
\[
\operatorname{tang}\beta\sin(\alpha''-L)+\operatorname{tang}B\sin(\alpha-\alpha'')+\operatorname{tang}\beta''\sin(L-\alpha).
\]
Simili modo characterem mutamus, si in expressione prima pro \emph{duabus} longitudinibus et latitudinibus geocentricis duae quaecunque heliocentricae terrae substituuntur. Si duae longitudines et latitudines in eandem expressionem ingredientes tantummodo inter se permutantur, etiam in charactere numeros respondentes permutare oportet: hinc autem valor ipse non mutatur, sed tantummodo e positivo negativus, e negativo positivus evadit. Ita e. g. fit $(0.1.2)=-(0.2.1)=(1.2.0)=-(1.0.2)=(2.0.1)=-(2.1.0)$. Omnes itaque quantitates hoc modo oriundae ad sequentes 19 reducuntur
\[
\begin{aligned}
&(0.1.2)\\
&(0.1.\mathrm{O}),\ (0.1.\mathrm{I}),\ (0.1.\mathrm{II}),\\
&(0.\mathrm{O}.2),\ (0.\mathrm{I}.2),\ (0.\mathrm{II}.2),\\
&(\mathrm{O}.1.2),\ (\mathrm{I}.1.2),\ (\mathrm{II}.1.2)\\
&(0.\mathrm{O}.\mathrm{I}),\ (0.\mathrm{O}.\mathrm{II}),\ (0.\mathrm{I}.\mathrm{II}),\\
&(1.\mathrm{O}.\mathrm{I}),\ (1.\mathrm{O}.\mathrm{II}),\ (1.\mathrm{I}.\mathrm{II}),\\
&(2.\mathrm{O}.\mathrm{I}),\ (2.\mathrm{O}.\mathrm{II}),\ (2.\mathrm{I}.\mathrm{II}),
\end{aligned}
\]
quibus accedit vigesima $(\mathrm{O}.\mathrm{I}.\mathrm{II})$.
\par
Ceterum facile demonstratur, singulas has expressiones, per productum e tribus cosinibus latitudinum ipsas ingredientium multiplicatas, aequales fieri volumini sextuplo pyramidis, cuius vertex est in Sole, basis vero triangulum formatum inter tria sphaerae coelestis puncta, quae locis expressionem illam ingredientibus respondent, statuto sphaerae radio $=1$. Quoties itaque hi tres \SourceLeaf{159}loci in eodem circulo maximo iacent, valor expressionis fieri debet $=0$; quod quum in tribus locis heliocentricis terrae semper locum habeat, quoties ad parallaxes et latitudines terrae a perturbationibus ortas non respicimus, i. e. quoties terram in ipso eclipticae plano constituimus, semper, hacce suppositione valente, erit $(\mathrm{O}.\mathrm{I}.\mathrm{II})=0$, quae quidem aequatio identica est, si pro plano tertio ecliptica ipsa accepta fuit. Ceterum quoties tum $B$, tum $B'$, tum $B''=0$, omnes istae expressiones, prima excepta, multo simpliciores fiunt; singulae scilicet a secunda usque ad decimam binis partibus conflatae erunt, ab undecima autem usque ad undevigesimam unico termino constabunt.
